\chapter{Infrastructure d'entraînement}\label{ch:ml:training-infrastructure}

Le modèle à carnet d'ordres d'une équipe s'entraîne pendant une journée, et l'équipe en chiffre un second \omterm{def:ml:training-infrastructure:accelerator}{accélérateur}. Avant d'en acheter un, elle doit savoir
où passe le temps. Sur la version réduite du problème du chapitre, mesurée sur un ordinateur portable, un tiers de chaque étape d'entraînement est
consacré à l'assemblage du lot, et non au calcul~; changer le format de stockage ne change presque rien, et rassembler le lot en une seule
lecture indexée plutôt que fenêtre par fenêtre ramène la part des données à un cinquième. Ce chapitre porte sur la mécanique qui se cache sous les
modèles des parties I à IV~: \omterm{def:ml:training-infrastructure:accelerator}{accélérateurs} et où passe le temps d'entraînement, \omterm{def:ml:training-infrastructure:loader}{chargeurs de données} pour les données tick qui ne laissent jamais une fenêtre franchir
une séance, l'arithmétique à précision réduite, l'entraînement sur plusieurs dispositifs, et des \omterm{def:ml:training-infrastructure:ckpt}{points de reprise d'entraînement} (checkpoints) qui permettent à une exécution interrompue de reprendre à l'identique.
Tout tourne sur un seul thread CPU, sans processus workers~; les parties multi-dispositifs sont simulées dans un seul processus, ce qui
suffit à vérifier leur arithmétique, et les temps sont ceux du portable de l'auteur, signalés comme tels.

\section{Accélérateurs et où passe le temps}

\begin{definition}[Accélérateur]\label{def:ml:training-infrastructure:accelerator}
Un \emph{accélérateur}\index{accélérateur} (accelerator) est un processeur conçu pour l'algèbre linéaire dense des réseaux de neurones (processeur graphique
ou processeur tensoriel dédié)~: des milliers d'unités arithmétiques, une mémoire à haute bande passante à leurs côtés, et
des unités matricielles à faible précision qui font l'essentiel du travail.
\end{definition}

\begin{dated}{2026-09}{Un accélérateur de centre de données}\label{dat:ml:training-infrastructure:h100}
La page de spécifications de NVIDIA pour le H100 indique 80~Go de mémoire avec 3.35~To/s de bande passante pour la version SXM (94~Go et
3.9~To/s pour la version NVL), et 1\,979 téraflops d'arithmétique \omterm{def:ml:training-infrastructure:mixed}{bfloat16} sur cœurs tenseurs pour la version SXM, un chiffre que la page
signale comme valeur avec sparsité.
\end{dated}

Des chiffres de ce genre expliquent le goulot d'étranglement habituel. Un \omterm{def:ml:training-infrastructure:accelerator}{accélérateur} peut consommer les données bien plus vite qu'un système de stockage, un
décompresseur ou une boucle Python peut l'alimenter, si bien que la première question face à un entraînement lent est la part de chaque étape pendant laquelle
le dispositif attend. L'étape se divise en chargement (lecture des données depuis le stockage), constitution des lots (découpe des fenêtres et leur stacking) et
calcul (passe avant, passe arrière, mise à jour de l'optimiseur). Le chapitre mesure les trois sur sa propre tâche~: huit
séances de dix minutes \texttt{firm.tape} du marché synthétique du livre, chacune une table de 1\,140 lignes d'une demi-seconde et de 41 colonnes (le
prix milieu et le carnet à dix niveaux du chapitre~8), et un réseau à deux couches qui lit les vingt dernières lignes et prédit la
variation du mid sur les dix lignes suivantes.

\section{Chargeurs de données pour les données tick}

\begin{definition}[Chargeur de données, jeu de données projeté en mémoire]\label{def:ml:training-infrastructure:loader}
Un \emph{chargeur de données}\index{chargeur de données} (data loader) est le code qui transforme les données stockées en lots d'entraînement~: il lit, décode, échantillonne,
découpe et empile les exemples, idéalement pendant que le dispositif calcule sur le lot précédent. Un \emph{jeu de données
projeté en mémoire}\index{jeu de données projeté en mémoire} (memory-mapped dataset) stocke les données dans un seul fichier binaire que le système d'exploitation projette dans l'espace
d'adressage du processus~: rien n'est lu tant qu'une tranche n'est pas touchée, le système met en cache ce qui est lu, et plusieurs processus peuvent
partager une seule copie.
\end{definition}

Les données tick imposent une règle au chargeur~: une fenêtre d'entraînement ne doit pas franchir la frontière d'une séance, sans quoi une fenêtre contiendra la
clôture d'une journée et l'ouverture de la suivante comme si elles étaient consécutives (et, avec des séances qui se chevauchent pour plusieurs instruments,
elle fuit d'un instrument à l'autre). L'échantillonneur présenté dans le \cref{lst:ml:train:sampler} tire
uniformément des couples (séance, début) parmi les fenêtres qui y tiennent, et son état aléatoire fait partie de ce qu'un \omterm{def:ml:training-infrastructure:ckpt}{point de reprise d'entraînement} enregistre.

Les séances sont stockées de trois façons~: un fichier CSV par séance (2.21~Mo au total), un fichier NumPy par séance (1.50~Mo), et un
fichier float32 projeté en mémoire, avec un petit index JSON (1.50~Mo). La \Cref{fig:ml:train:time} montre le coût mesuré d'une étape d'entraînement
avec chacun, le chargement étant compté une fois par \omterm{def:ml:neural-networks-for-noisy-tabular-data:backprop}{époque} de 100 étapes. L'analyse syntaxique des fichiers CSV prend 12.5~ms par \omterm{def:ml:neural-networks-for-noisy-tabular-data:backprop}{époque} (176~Mo/s),
la lecture des fichiers NumPy prend 0.83~ms (1.8~Go/s), et la projection du fichier unique 0.39~ms quelle que soit sa taille. Par étape, ces valeurs sont minimes
face à la demi-milliseconde nécessaire pour découper 256 fenêtres une à une en Python et aux 1.1 à 1.2~ms de calcul~: la part des données dans
l'étape est de 0.34, 0.31 et 0.39. Rassembler le lot entier depuis le fichier projeté en une seule lecture indexée
(\cref{lst:ml:train:gather}) prend 0.27~ms et ramène la part à 0.20.

\begin{omfigure}
\begin{tikzpicture}
\begin{axis}[width=11cm,height=6cm,ybar stacked,bar width=22pt,ymin=0,ymax=2.6,ylabel={millisecondes par étape d'entraînement},
  area legend,xtick={0,1,2,3},xticklabels={{CSV},{NumPy},{projection\\en mémoire},{projection\\en mémoire,\\lecture groupée}},xticklabel style={align=center},
  tick label style={font=\small},label style={font=\small},xmin=-0.6,xmax=3.6,grid=major,grid style={gray!25},
  legend columns=3,legend style={font=\footnotesize,at={(0.5,-0.3)},anchor=north,draw=none,fill=none,/tikz/every even column/.append style={column sep=8pt}},table/col sep=comma]
\addplot[fill=omThm,draw=none] table[x=i,y expr=\thisrow{load_ms}/100]{figdata/ml/23-training-infrastructure/measured_loader.csv};
\addplot[fill=omProp,draw=none] table[x=i,y=batch_ms]{figdata/ml/23-training-infrastructure/measured_loader.csv};
\addplot[fill=omDef,draw=none] table[x=i,y=compute_ms]{figdata/ml/23-training-infrastructure/measured_loader.csv};
\legend{chargement (par étape),constitution des lots,calcul}
\end{axis}
\end{tikzpicture}

\omcaption{Où passe le temps d'une étape d'entraînement, selon le format de stockage et la constitution des lots~: médianes de cinq \omterm{def:ml:neural-networks-for-noisy-tabular-data:backprop}{époques} de 100 étapes de 256
fenêtres. Mesures faites sur un ordinateur portable (Intel Core Ultra 7 155H) sous WSL2, un seul thread, sans cœurs isolés~; d'autres tâches ont pu tourner.
Données~: \texttt{bench\_train.py}, \texttt{measured\_loader.csv}.}\label{fig:ml:train:time}
\end{omfigure}

À cette échelle, la leçon porte sur le code et non sur le stockage~: la boucle Python sur les fenêtres coûte plus cher que n'importe quel format. À l'échelle d'une
année de données de carnet d'ordres, les formats divergent. Une année de la même table au pas d'une demi-seconde pour un instrument, 250 séances de six heures et
demie, est d'environ 1.9~Go en float32 et 2.8~Go en texte~: aux débits mesurés, l'analyse du texte coûte environ 16~secondes par
\omterm{def:ml:neural-networks-for-noisy-tabular-data:backprop}{époque}, la lecture des fichiers binaires prend environ une seconde, et la projection ne coûte rien tant que les fenêtres ne sont pas lues, à partir du cache de pages après la première \omterm{def:ml:neural-networks-for-noisy-tabular-data:backprop}{époque}.
Et au-delà de la mémoire de la machine, seul un format projeté en mémoire ou diffusé en flux continu fonctionne tout simplement.

\section{Précision mixte}

\begin{definition}[Entraînement en précision mixte, bfloat16]\label{def:ml:training-infrastructure:mixed}
L'\emph{entraînement en précision mixte}\index{entraînement en précision mixte} (mixed-precision training) exécute les opérations coûteuses (produits matriciels, convolutions)
dans un format de 16 bits tout en conservant les poids, la perte et l'état de l'optimiseur en 32 bits (Micikevicius et coauteurs,
2018). \emph{bfloat16}\index{bfloat16} est un format flottant sur 16 bits, avec les huit bits d'exposant de float32 et sept bits
de mantisse~: la même plage que float32, et un epsilon machine (livre~4, chapitre~25) de $2^{-7}\approx0.0078$ au lieu de $2^{-23}$, si bien qu'aucune
mise à l'échelle de la perte n'est requise pour prévenir le sous-dépassement (Kalamkar et coauteurs, 2019).
\end{definition}

Sous l'autocast \omterm{def:ml:training-infrastructure:mixed}{bfloat16} du processeur, le réseau du chapitre s'entraîne au même point (la \cref{fig:ml:train:bf16})~: sur les 50 dernières
étapes sur 300, la perte moyenne vaut 0.0525 contre 0.0527 en float32. Pas à pas, les deux exécutions diffèrent~: le plus grand écart
entre leurs pertes sur l'exécution est de 0.0132, soit 22~\% de la perte float32 de cette étape, car l'arrondi modifie légèrement chaque mise à jour et les
trajectoires s'écartent. Reproduire une exécution bit à bit d'une précision à l'autre n'est pas à espérer~; l'équivalence s'apprécie sur la
métrique de validation, sur toute l'exécution.

\begin{omfigure}
\begin{tikzpicture}
\begin{axis}[width=11cm,height=5.6cm,xlabel={étape d'entraînement},ylabel={perte (moyenne sur 10 étapes)},tick label style={font=\small},
  label style={font=\small},xmin=10,xmax=300,grid=major,grid style={gray!25},legend cell align=left,
  legend style={font=\footnotesize,at={(0.98,0.98)},anchor=north east,fill=white,draw=none},table/col sep=comma]
\addplot[omDef,very thick] table[x=step,y=fp32]{figdata/ml/23-training-infrastructure/losses.csv};
\addplot[omThm,thick,dashed] table[x=step,y=bf16]{figdata/ml/23-training-infrastructure/losses.csv};
\legend{float32,bfloat16 sous autocast}
\end{axis}
\end{tikzpicture}

\omcaption{Perte d'entraînement du même réseau, du même ordre des données et de la même initialisation, en float32 et sous autocast \omterm{def:ml:training-infrastructure:mixed}{bfloat16}.
Données~: \texttt{ml\_train.run} (\texttt{fig\_train.py}).}\label{fig:ml:train:bf16}
\end{omfigure}

\section{Entraînement distribué}

\begin{definition}[Accumulation de gradient, parallélisme de données, parallélisme de modèle, all-reduce]\label{def:ml:training-infrastructure:dist}
L'\emph{accumulation de gradient}\index{accumulation de gradient} somme les gradients de plusieurs micro-lots avant une seule étape de l'optimiseur,
ce qui équivaut à une étape sur leur union. Le \emph{parallélisme de données}\index{parallélisme de données} fait tourner une copie du modèle sur chaque dispositif,
chacune sur sa propre part du lot, et moyenne leurs gradients avant chaque étape (Li et coauteurs, 2020)~; le \emph{parallélisme de
modèle}\index{parallélisme de modèle} répartit les couches ou les matrices d'un même modèle entre plusieurs dispositifs lorsqu'il ne tient pas sur un seul
(Shoeybi et coauteurs, 2019). Un \emph{all-reduce}\index{all-reduce} est l'opération collective qui laisse chaque dispositif avec
la somme (ou la moyenne) de tous les tableaux des dispositifs.
\end{definition}

Les deux identités sont vérifiées dans le code du chapitre (\cref{lst:ml:train:loop}). Quatre micro-lots accumulés de 64 et un lot
de 256 des mêmes fenêtres laissent des paramètres qui diffèrent d'au plus $1.4\times10^{-7}$ après 20 étapes, l'arrondi intervenant
dans un ordre de sommation différent. La moyenne des gradients de deux travailleurs sur les moitiés d'un lot diffère du gradient du lot
entier d'au plus $1.5\times10^{-8}$, pour des gradients de taille 0.12~: un \omterm{def:ml:training-infrastructure:dist}{all-reduce} calcule exactement le gradient du grand lot, et le
\omterm{def:ml:training-infrastructure:dist}{parallélisme de données} est une \omterm{def:ml:training-infrastructure:dist}{accumulation de gradient} répartie sur plusieurs dispositifs. Ce qui change avec un lot effectif plus grand,
c'est l'optimisation~: le \omterm{def:ml:trees-and-boosting:boosting}{taux d'apprentissage} et sa phase de montée doivent être remis à l'échelle (Goyal et coauteurs, 2017), et pour les
petits modèles bruités du trading, des lots plus grands aident rarement.

\section{Points de reprise et reprise déterministe}

\begin{definition}[Point de reprise d'entraînement]\label{def:ml:training-infrastructure:ckpt}
Un \emph{point de reprise d'entraînement}\index{point de reprise d'entraînement} est un état sauvegardé à partir duquel un calcul peut reprendre comme s'il ne s'était jamais arrêté~:
les paramètres du modèle, l'état de l'optimiseur (estimations des moments, nombres de pas), la position et l'état aléatoire de l'échantillonneur de données,
les états des générateurs aléatoires et le nombre de pas.
\end{definition}

Un calcul de 200 étapes et un calcul interrompu à l'étape 100, sauvegardé en point de reprise (\cref{lst:ml:train:ckpt}), puis repris dans des
objets nouvellement construits avec une initialisation différente, se terminent avec des paramètres identiques bit à bit et des pertes identiques
étape par étape~; le point de reprise pèse 693~ko. Ne sauvegarder que les poids du modèle, le raccourci habituel, réinitialise les estimations
des moments d'\omterm{def:ml:neural-networks-for-noisy-tabular-data:adam}{Adam} et rejoue les premières fenêtres de l'échantillonneur~; le calcul repris finit avec des paramètres éloignés d'au plus 0.039 de
ceux du calcul non interrompu. Pour la recherche qui doit être reproduite (chapitre~25) et pour les travaux sur des machines qui peuvent être
préemptées, le point de reprise complet fait la différence entre reprendre et recommencer avec un calcul différent.

\begin{method}[Avant d'acheter un accélérateur]\label{met:ml:train:method}
\begin{enumerate}
\item Mesurer la charge, la taille des lots et le temps de calcul par étape~; si les données prennent plus qu'une petite part, corriger d'abord le
chargeur (stockage binaire ou projeté en mémoire, regroupement des lots en une seule lecture, préchargement).
\item Essayer l'autocast \omterm{def:ml:training-infrastructure:mixed}{bfloat16} et le juger sur la métrique de validation au fil du calcul.
\item Agrandir le lot par accumulation avant d'agrandir le matériel~; remettre à l'échelle le \omterm{def:ml:trees-and-boosting:boosting}{taux d'apprentissage} et vérifier que le modèle s'améliore.
\item Sauvegarder tout ce qui détermine l'avenir du calcul, et tester la reprise bit à bit.
\end{enumerate}
\end{method}

\section{Tutoriel~: où part le temps}\label{tut:ml:training-infrastructure:tut}

\begin{tutorial}
\textbf{Objectif.} Stocker les sessions de cotations de trois façons, mesurer le budget temps d'une étape, entraîner en \omterm{def:ml:training-infrastructure:mixed}{bfloat16}, vérifier l'accumulation et le
\omterm{def:ml:training-infrastructure:dist}{parallélisme de données}, et reprendre un calcul bit à bit. \textbf{État final~:} \cref{fig:ml:train:time,fig:ml:train:bf16}.
\begin{tutsteps}
\item \textbf{Un échantillonneur de fenêtres conscient des sessions, avec un état sauvegardable.}
\omcode{code/firm/mltrain/firm_mltrain.py}{78}{96}{Fenêtres qui ne franchissent jamais une frontière de session.}\label{lst:ml:train:sampler}
\item \textbf{Une seule collecte pour un lot entier.}
\omcode{code/firm/mltrain/firm_mltrain.py}{110}{116}{Un lot tiré du jeu de données projeté en mémoire en une seule lecture indexée.}\label{lst:ml:train:gather}
\item \textbf{Accumulation et autocast \omterm{def:ml:training-infrastructure:mixed}{bfloat16}.}
\omcode{code/firm/mltrain/firm_mltrain.py}{119}{141}{La boucle d'entraînement.}\label{lst:ml:train:loop}
\item \textbf{Points de reprise.}
\omcode{code/firm/mltrain/firm_mltrain.py}{156}{167}{Sauvegarde et restauration de tout ce dont dépend l'avenir du calcul.}\label{lst:ml:train:ckpt}
\item \textbf{Exécuter} \texttt{ml\_train.precision()}, \texttt{accumulation()}, \texttt{data\_parallel()}, \texttt{resume()},
\texttt{naive\_resume()}, \texttt{fig\_train.py}~; et, une fois, \texttt{bench\_train.py} sous \texttt{nice}.
\end{tutsteps}
\textbf{Pour aller plus loin.} Précharger le lot suivant dans un thread d'arrière-plan et mesurer le recouvrement~; ajouter un second instrument
et vérifier que les fenêtres ne mélangent jamais les instruments.
\end{tutorial}

\section{Construction~: la boîte à outils d'entraînement}\label{bld:ml:training-infrastructure:mltrain}

\begin{build}
\bfield{Rôle} Des calculs d'entraînement rapides là où cela compte, corrects aux frontières de session, et reprenables exactement.
\bfield{Interface} \texttt{write\_shards(sessions, root, fmt)}, \texttt{Shards(root, fmt)},
\texttt{WindowSampler(lengths, T, horizon, batch, seed)} avec \texttt{state}/\texttt{set\_state}, \texttt{make\_batch}, \texttt{make\_batch\_vec},
\texttt{train(model, shards, sampler, steps, lr, accumulate, bf16)}, \texttt{data\_parallel\_grads}, \texttt{save\_checkpoint},
\texttt{load\_checkpoint}.
\bfield{Règles} Les fenêtres ne franchissent jamais les sessions~; les durées sont mesurées, étiquetées avec la machine, et jamais affirmées~; tout ce qui
détermine la suite d'un calcul entre dans son point de reprise.
\bfield{Tests d'acceptation} \texttt{code/firm/mltrain/tests/}~: les trois formats renvoient des fenêtres identiques~; aucune fenêtre échantillonnée ne franchit
une frontière~; le lot collecté égale le lot bouclé~; l'accumulation égale un lot plus grand~; les gradients moyens des fragments égalent le
gradient complet~; la reprise est bit à bit.
\bfield{Pour aller plus loin} \omterm{def:ml:training-infrastructure:dist}{Parallélisme de données} réel à deux processus avec le backend gloo lorsque la machine le permet~; diffusion en flux depuis des
fragments compressés~; préchargement.
\end{build}

\begin{omsources}
\item P.~Micikevicius et coauteurs, ``Mixed precision training'', arXiv:1710.03740, 2018.
\item D.~Kalamkar et coauteurs, ``A study of BFLOAT16 for deep learning training'', arXiv:1905.12322, 2019.
\item S.~Li et coauteurs, ``PyTorch distributed: experiences on accelerating data parallel training'', arXiv:2006.15704, 2020.
\item P.~Goyal et coauteurs, ``Accurate, large minibatch SGD: training ImageNet in 1 hour'', arXiv:1706.02677, 2017.
\item M.~Shoeybi et coauteurs, ``Megatron-LM: training multi-billion parameter language models using model parallelism'',
arXiv:1909.08053, 2019.
\item NVIDIA, H100 Tensor Core GPU specifications (web page, accessed September 2026).
\end{omsources}

\section{Exercices}

\begin{exercise}[$\star$]\label{exo:ml:training-infrastructure:1}
Quel est l'epsilon machine du \omterm{def:ml:training-infrastructure:mixed}{bfloat16}, du float16 et du float32~? Lequel d'entre eux peut-il représenter $10^{-30}$~?
\end{exercise}

\begin{exercise}[$\star$]\label{exo:ml:training-infrastructure:2}
Une session compte 1\,140 lignes~; les fenêtres font 20 lignes, avec un horizon cible de 10 lignes. Combien de fenêtres contient-elle~? Que se passe-t-il si
un échantillonneur tire les débuts uniformément sur l'ensemble des sessions concaténées~?
\end{exercise}

\begin{exercise}[$\star$]\label{exo:ml:training-infrastructure:3}
D'après la \cref{fig:ml:train:time}, quelle accélération donnerait la division par deux du temps de calcul avec le chargeur bouclé, et avec la collecte~?
\end{exercise}

\begin{exercise}[$\star\star$]\label{exo:ml:training-infrastructure:4}
Pourquoi les pertes en float32 et en \omterm{def:ml:training-infrastructure:mixed}{bfloat16} diffèrent-elles jusqu'à 22~\% à certaines étapes alors que leurs moyennes concordent~?
\end{exercise}

\begin{exercise}[$\star\star$]\label{exo:ml:training-infrastructure:5}
Montrer que, avec une perte moyenne, accumuler $k$ micro-lots pondérés chacun par $1/k$ donne le gradient de la moyenne sur leur union.
Qu'est-ce qui change avec la \omterm{def:ml:neural-networks-for-noisy-tabular-data:dropout}{normalisation par lot} (batch normalisation)~?
\end{exercise}

\begin{exercise}[$\star\star$]\label{exo:ml:training-infrastructure:6}
\emph{Trouver la faille.} «~Nous reprenons les tâches préemptées à partir des derniers poids sauvegardés, de sorte que nos exécutions sont reproductibles.~»
\end{exercise}

\begin{exercise}[$\star\star\star$]\label{exo:ml:training-infrastructure:7}
Estimer, à partir des débits mesurés, le temps de chargement par \omterm{def:ml:neural-networks-for-noisy-tabular-data:backprop}{époque} d'une année du tableau du chapitre dans chaque format (250
sessions de six heures et demie de lignes d'une demi-seconde, 41 colonnes en float32, texte environ 1.5 fois plus volumineux que le binaire).
\end{exercise}

\begin{exercise}[$\star\star\star$]\label{exo:ml:training-infrastructure:8}
Dans un \omterm{def:ml:training-infrastructure:dist}{all-reduce} en anneau sur $p$ dispositifs détenant chacun $n$ nombres, chaque dispositif envoie et reçoit environ $2n(p-1)/p$ nombres. Expliquer
pourquoi, et ce que cela implique pour la manière dont l'entraînement en \omterm{def:ml:training-infrastructure:dist}{parallélisme de données} passe à l'échelle avec le nombre de dispositifs.
\end{exercise}

\section{Problème~: Où passe le temps}

\begin{problem}[{Problème du week-end --- mesurer avant d'acheter}]\label{pb:ml:training-infrastructure:1}
Les sessions, les formats, le réseau et les exécutions du chapitre.

\medskip\noindent\textbf{Partie I --- Le budget.}
\begin{enumerate}
\item Quelle est la taille des données dans chaque format~?
\item Quels sont les temps de chargement, de traitement par lot et de calcul par étape~?
\item Quelle part de l'étape est consacrée aux données pour chaque format de stockage, et avec le gather~?
\item Que faudrait-il changer en premier, et pourquoi~?
\end{enumerate}

\medskip\noindent\textbf{Partie II --- Précision.}
\begin{enumerate}[resume]
\item Quels sont les bits, la plage et l'epsilon de \omterm{def:ml:training-infrastructure:mixed}{bfloat16}~?
\item Comment se comparent les exécutions en float32 et en \omterm{def:ml:training-infrastructure:mixed}{bfloat16}~?
\item Pourquoi aucune mise à l'échelle de la perte (loss scaling) n'est-elle nécessaire avec \omterm{def:ml:training-infrastructure:mixed}{bfloat16}~?
\item Où conserverait-on le float32 dans le pipeline d'un modèle de trading~?
\end{enumerate}

\medskip\noindent\textbf{Partie III --- Parallélisme et points de reprise.}
\begin{enumerate}[resume]
\item Que montrent les contrôles d'accumulation et de \omterm{def:ml:training-infrastructure:dist}{parallélisme de données}~?
\item Que doit contenir un \omterm{def:ml:training-infrastructure:ckpt}{point de reprise d'entraînement} (checkpoint), et que coûte l'omission d'une partie de celui-ci~?
\item Quelle est la taille du point de reprise du chapitre~?
\item Qu'exige un lot effectif plus grand du \omterm{def:ml:trees-and-boosting:boosting}{taux d'apprentissage}~?
\end{enumerate}

\medskip\noindent\textbf{Partie IV --- Le verdict.}
\begin{enumerate}[resume]
\item Énoncer le \emph{résultat nommé}~: la part du chargement des données dans le temps d'étape pour chaque format de stockage, et le plus grand écart de perte
entre les entraînements en float32 et en \omterm{def:ml:training-infrastructure:mixed}{bfloat16} sur l'ensemble de l'exécution.
\item L'équipe aurait-elle besoin d'un second \omterm{def:ml:training-infrastructure:accelerator}{accélérateur}~?
\item Quand le \omterm{def:ml:training-infrastructure:dist}{parallélisme de modèle} devient-il nécessaire pour un modèle de trading~?
\item Comment rendre reproductible, d'une machine à l'autre, une exécution d'entraînement~?
\item Quel est le risque d'un échantillonneur qui chevauche des sessions~?
\item Comment mesurer un \omterm{def:ml:training-infrastructure:loader}{chargeur de données} (data loader) sur une infrastructure partagée~?
\item Que faudrait-il journaliser pour chaque exécution d'entraînement~?
\item En une phrase~: quelle est la première chose à mesurer sur une tâche d'entraînement lente~?
\end{enumerate}
\end{problem}

\section{Questions d'entretien}

\begin{interviewq}[$\star$ \iqroles{ingénieur ML}]\label{iq:ml:training-infrastructure:1}
Qu'est-ce que \omterm{def:ml:training-infrastructure:mixed}{bfloat16} et pourquoi lui préfère-t-on float16 pour l'entraînement~?
\end{interviewq}

\begin{interviewq}[$\star\star$ \iqroles{ingénieur ML, développeur}]\label{iq:ml:training-infrastructure:2}
Votre taux d'utilisation des GPU est de 30~\%. Comment en trouvez-vous la cause~?
\end{interviewq}

\begin{interviewq}[$\star\star$ \iqroles{ingénieur ML}]\label{iq:ml:training-infrastructure:3}
Expliquez le \omterm{def:ml:training-infrastructure:dist}{parallélisme de données} et ce que fait un \omterm{def:ml:training-infrastructure:dist}{all-reduce}.
\end{interviewq}

\begin{interviewq}[$\star\star$ \iqroles{ingénieur ML, chercheur}]\label{iq:ml:training-infrastructure:4}
Comment stockeriez-vous et échantillonneriez-vous une année de données de carnet d'ordres pour entraîner des modèles de séquences~?
\end{interviewq}

\begin{interviewq}[$\star\star$ \iqroles{ingénieur ML}]\label{iq:ml:training-infrastructure:5}
Que doit contenir un \omterm{def:ml:training-infrastructure:ckpt}{point de reprise d'entraînement} (checkpoint) pour une reprise bit à bit~?
\end{interviewq}

\begin{interviewq}[$\star\star\star$ \iqroles{ingénieur ML, développeur}]\label{iq:ml:training-infrastructure:6}
On considère souvent que l'\omterm{def:ml:training-infrastructure:dist}{accumulation de gradient} et le \omterm{def:ml:training-infrastructure:dist}{parallélisme de données} sont équivalents. Dans quels cas ne le sont-ils pas~?
\end{interviewq}
